\documentclass[12pt,a4paper,oneside]{report}

% Informe técnico de despliegue de VehicleSense en AWS EC2.
\usepackage[utf8]{inputenc}
\usepackage[T1]{fontenc}
\usepackage[spanish,es-nodecimaldot]{babel}
\usepackage{lmodern}
\usepackage{microtype}
\usepackage[a4paper,margin=2.6cm,headheight=15pt]{geometry}
\usepackage{graphicx}
\usepackage{amsmath}
\usepackage{booktabs}
\usepackage{tabularx}
\usepackage{longtable}
\usepackage{array}
\usepackage{enumitem}
\usepackage{xcolor}
\usepackage{listings}
\usepackage{fancyhdr}
\usepackage{lastpage}
\usepackage[hidelinks]{hyperref}

\definecolor{awsorange}{HTML}{D97706}
\definecolor{navy}{HTML}{16324F}
\definecolor{ink}{HTML}{263238}
\definecolor{softgray}{HTML}{F3F5F7}
\definecolor{linegray}{HTML}{D9E0E6}
\definecolor{codegreen}{HTML}{176B4D}

\hypersetup{
  colorlinks=true,
  linkcolor=navy,
  citecolor=navy,
  urlcolor=navy,
  pdftitle={Despliegue de VehicleSense en AWS EC2},
  pdfauthor={Patricio Vasquez}
}

\setlength{\parindent}{0pt}
\setlength{\parskip}{0.55em}
\setlist[itemize]{leftmargin=1.4em,itemsep=0.25em,topsep=0.25em}
\setlist[enumerate]{leftmargin=1.7em,itemsep=0.35em,topsep=0.25em}
\renewcommand{\arraystretch}{1.2}

\pagestyle{fancy}
\fancyhf{}
\fancyhead[L]{\small\textcolor{navy}{VehicleSense}}
\fancyhead[R]{\small\textcolor{navy}{Despliegue en AWS EC2}}
\fancyfoot[C]{\small\thepage\ de \pageref{LastPage}}
\renewcommand{\headrulewidth}{0.4pt}
\renewcommand{\footrulewidth}{0pt}

\lstdefinestyle{terminal}{
  backgroundcolor=\color{softgray},
  basicstyle=\ttfamily\small,
  breaklines=true,
  columns=fullflexible,
  frame=single,
  rulecolor=\color{linegray},
  keepspaces=true,
  showstringspaces=false,
  numbers=none,
  xleftmargin=0.4em,
  xrightmargin=0.4em
}
\lstset{style=terminal}

\newcommand{\status}[1]{\textcolor{codegreen}{\textbf{#1}}}
\newcommand{\important}[1]{\textcolor{awsorange}{\textbf{#1}}}
\newcommand{\code}[1]{\texttt{\small #1}}

\begin{document}

\begin{titlepage}
  \thispagestyle{empty}
  \vspace*{1.4cm}
  \begin{center}
    {\Large\textcolor{navy}{PROYECTO DE MONITOREO VEHICULAR}}\\[1.2cm]
    {\Huge\bfseries\textcolor{navy}{Despliegue de la plataforma web\\
    VehicleSense en AWS EC2}\par}
    \vspace{0.7cm}
    {\large Informe técnico final sobre arquitectura cloud, operación y validación}\par
    \vfill
    \rule{0.72\textwidth}{0.8pt}\par
    \vspace{0.35cm}
    {\large Patricio Vasquez}\par
    \vspace{0.25cm}
    {\large 31 de agosto de 2026}
  \end{center}
\end{titlepage}

\pagenumbering{roman}
\chapter*{Resumen}
\addcontentsline{toc}{chapter}{Resumen}

Este informe documenta el despliegue productivo de VehicleSense, una plataforma de
monitoreo vehicular que integra un dispositivo ESP32, sensores, transporte MQTT,
un backend FastAPI, PostgreSQL y una interfaz React. La aplicación fue publicada
en una instancia AWS EC2 Ubuntu 24.04 ARM64 mediante Docker Compose. Nginx actúa
como punto único de entrada y entrega el frontend, la API REST y el canal
WebSocket; HiveMQ Cloud permanece como broker MQTT externo.

Se describen la arquitectura, la configuración de red, el uso de Elastic IP,
Security Groups, DNS de Namecheap, certificados Let's Encrypt, migraciones
Alembic, backups y operación diaria. La elección de AWS EC2 se fundamenta en el
control sobre el sistema Linux y la red, la compatibilidad con ARM64, la
integración de sus servicios y el valor formativo de administrar directamente el
servidor. La validación confirmó el acceso HTTPS, los endpoints de salud, la
conexión del backend con PostgreSQL y HiveMQ, y la independencia de la Mac local.

El despliegue queda operativo como una topología de una sola instancia. Por ello,
la disponibilidad observada no equivale a alta disponibilidad y la precisión de
los sensores, la ejecución de comandos del firmware y la validación física del
dispositivo permanecen fuera del alcance de este informe.

\textbf{Palabras clave:} AWS EC2, Docker Compose, Nginx, FastAPI, MQTT, HiveMQ,
PostgreSQL, React, HTTPS, monitoreo vehicular.

\tableofcontents
\clearpage
\pagenumbering{arabic}

\chapter{Introducción}

\section{Planteamiento}

VehicleSense nació como un prototipo académico para observar variables de un
vehículo y presentarlas en una interfaz web. El sistema combina mediciones de
temperatura, humedad, posición, aceleración, iluminación y características
acústicas. Para que la información pueda consultarse sin mantener un computador
personal encendido, fue necesario trasladar la capa cloud a una infraestructura
conectada permanentemente a Internet.

El reto no consistió únicamente en copiar el frontend a un servidor. La solución
debía recibir mensajes MQTT sobre TLS, validar los contratos de telemetría,
persistir el historial, autenticar usuarios, servir una API y mantener una conexión
WebSocket para actualizaciones en vivo. Además, el acceso público debía utilizar
HTTPS y no debía exponer PostgreSQL, FastAPI ni el frontend interno directamente.

\section{Objetivos}

\begin{enumerate}
  \item Publicar la aplicación web en el dominio
        \url{https://vehiclemonitorsense.me} con HTTPS.
  \item Ejecutar frontend, backend y PostgreSQL en una instancia EC2 reproducible
        mediante Docker Compose.
  \item Conectar FastAPI con HiveMQ Cloud por MQTT/TLS en el puerto 8883.
  \item Proteger la superficie pública con Security Groups, Nginx, TLS, roles y
        separación de credenciales.
  \item Validar el funcionamiento mediante health checks, migraciones,
        autenticación, pruebas de red y procedimientos de backup.
\end{enumerate}

\section{Alcance y criterio de evidencia}

Este documento cubre el despliegue de la plataforma web y su operación inicial.
Cuando se indica que una característica está implementada, existe configuración o
código versionado. Cuando se indica que está validada en producción, existe una
comprobación fechada en la infraestructura desplegada. La validación de precisión
de sensores, el bring-up eléctrico, la calibración acústica y la ejecución
semántica de comandos no se presentan como resultados del despliegue cloud.

\chapter{Contexto técnico del proyecto}

\section{Flujo de información}

El flujo completo separa el dispositivo de la aplicación web. El ESP32 adquiere y
valida muestras; la telemetría se serializa según los esquemas del proyecto y se
publica mediante MQTT/TLS. FastAPI recibe los mensajes desde HiveMQ, verifica
identidades, tópicos, tamaño, esquema y valores finitos, y luego persiste los
datos en PostgreSQL. React consume REST y WebSocket a través de Nginx.

\begin{center}
\fbox{\begin{minipage}{0.92\textwidth}
\centering
\code{ESP32}\quad $\xrightarrow{\text{MQTT/TLS 8883}}$ \quad
\code{HiveMQ Cloud}\quad $\xrightarrow{\text{ingesta}}$ \quad
\code{FastAPI}\quad $\xrightarrow{\text{SQL}}$ \quad
\code{PostgreSQL}\\[0.45cm]
\code{Nginx}\quad $\xrightarrow{/}$ \quad \code{React} \qquad
\code{Nginx}\quad $\xrightarrow{/api, /ws}$ \quad \code{FastAPI}
\end{minipage}}
\end{center}

Los tópicos canónicos utilizados por el backend son:

\begin{lstlisting}
vehiclesense/v1/vehicles/{vehicle_id}/devices/{device_id}/telemetry
vehiclesense/v1/vehicles/{vehicle_id}/devices/{device_id}/status
vehiclesense/v1/vehicles/{vehicle_id}/devices/{device_id}/events
vehiclesense/v1/vehicles/{vehicle_id}/devices/{device_id}/acoustic
vehiclesense/v1/vehicles/{vehicle_id}/devices/{device_id}/commands
vehiclesense/v1/vehicles/{vehicle_id}/devices/{device_id}/command-acks
\end{lstlisting}

\section{Componentes desplegados}

\begin{table}[ht]
\centering
\small
\begin{tabularx}{\textwidth}{>{\bfseries}l X X}
\toprule
Componente & Función en el proyecto & Ubicación o tecnología \\
\midrule
ESP32 & Adquisición, validación, almacenamiento local y publicación de datos &
  Hardware del prototipo \\
HiveMQ Cloud & Autenticación, ACL, TLS y transporte MQTT & Servicio externo \\
FastAPI & Ingesta, autenticación, REST, WebSocket, alertas y comandos & Contenedor Docker \\
PostgreSQL & Persistencia de usuarios, dispositivos, telemetría, alertas y viajes & Contenedor Docker \\
React/Vite & Dashboard, mapas, gráficas, alertas y reportes & Contenedor estático Nginx \\
Nginx & Entrada HTTP/HTTPS, proxy REST/WS y redirecciones & Contenedor público \\
\bottomrule
\end{tabularx}
\caption{Componentes de la solución desplegada.}
\end{table}

\chapter{Arquitectura en AWS}

\section{Recursos de la instancia}

La plataforma se desplegó en una instancia dedicada de producción, con Ubuntu
Server 24.04 LTS sobre arquitectura
ARM64. Se seleccionó una instancia \code{t4g.small}, con 2 vCPU y 2 GiB de RAM,
y un volumen raíz gp3 de 30 GiB. El diseño no fija una arquitectura en Docker,
por lo que también es compatible con EC2 AMD64.

La instancia tiene una Elastic IP asociada. Esto permite que los registros DNS
continúen apuntando a una dirección estable aunque la instancia se detenga o se
reemplace. La Mac del desarrollador no participa en la ejecución: únicamente se
utiliza para administrar el servidor por SSH cuando es necesario.

\section{Red y exposición de puertos}

Se creó un Security Group dedicado. Las reglas de entrada se limitaron a lo
estrictamente necesario:

\begin{table}[ht]
\centering
\small
\begin{tabularx}{0.9\textwidth}{l c X}
\toprule
Servicio & Puerto & Origen o propósito \\
\midrule
SSH & 22/TCP & Solo la IP pública actual del administrador, en formato /32 \\
HTTP & 80/TCP & Público para redirección y desafío ACME \\
HTTPS & 443/TCP & Público para la aplicación web \\
\midrule
PostgreSQL & 5432/TCP & No publicado; red Docker interna \\
FastAPI & 8000/TCP & No publicado; solo Nginx y la red Docker \\
Frontend & 8080/TCP & No publicado; solo Nginx \\
MQTT & 8883/TCP & No es entrada EC2; conexión saliente a HiveMQ \\
\bottomrule
\end{tabularx}
\caption{Puertos y frontera de red de EC2.}
\end{table}

La red Docker \code{database} se define como interna. FastAPI participa tanto en
la red privada de base de datos como en la red de borde, mientras que PostgreSQL
solo participa en la red privada. Nginx es el único servicio que publica puertos
del host.

\section{DNS y dominio}

En Namecheap se reemplazaron los registros de parking por registros A para el
dominio raíz y para \code{www}. Ambos registros apuntan a la Elastic IP de la
instancia. La propagación se comprobó con \code{dig} antes de solicitar el
certificado. El comportamiento final es:

\begin{itemize}
  \item \code{http://vehiclemonitorsense.me} redirige a HTTPS.
  \item \code{http://www.vehiclemonitorsense.me} redirige a HTTPS.
  \item \code{https://www.vehiclemonitorsense.me} redirige al dominio principal.
\end{itemize}

\section{Persistencia y disponibilidad}

PostgreSQL se ejecuta con un volumen Docker persistente denominado
\code{postgres\_data}. Alembic aplica las migraciones antes de iniciar FastAPI; si
la migración falla, el backend no inicia. Los backups se generan con
\code{pg\_dump} en formato custom, se acompañan de un checksum SHA-256 y deben
copiarse a un almacenamiento cifrado externo.

La topología es de una sola instancia. Esto simplifica el proyecto académico y
reduce la cantidad de componentes, pero introduce un punto único de fallo. No se
afirma alta disponibilidad, balanceo ni recuperación automática entre zonas.

\chapter{Proceso de despliegue y manejo de AWS}

\section{Preparación de Ubuntu}

Después de lanzar la instancia y asociar la Elastic IP, se accedió por SSH con la
llave privada protegida mediante permisos \code{600}. En el servidor se creó
\code{/opt/vehiclesense}, se clonó el repositorio y se ejecutó el script versionado
\code{bootstrap\_ubuntu.sh}. El script instala Docker Engine, Buildx, Compose,
Git, OpenSSL, jq y los certificados base, y agrega el usuario \code{ubuntu} al
grupo \code{docker}.

La sesión SSH se cerró y se abrió nuevamente para que el grupo se aplicara. Las
versiones observadas en producción fueron Docker Engine 29.7.2 y Docker Compose
v5.5.0 sobre \code{linux/arm64}.

\section{Configuración de variables}

El archivo \code{/opt/vehiclesense/deploy/.env} no se versiona y se protegió con
\code{chmod 600}. Se configuraron:

\begin{itemize}
  \item dominio principal, dominio alternativo y correo de Let's Encrypt;
  \item credenciales y nombre de la base de datos;
  \item secreto JWT y usuario administrador inicial;
  \item host, puerto, usuario, contraseña y client ID de HiveMQ;
  \item orígenes CORS HTTPS y umbrales de estado de dispositivos.
\end{itemize}

Las credenciales del backend son diferentes de las credenciales compiladas en el
ESP32. No se copiaron secretos, certificados ni dumps a Git, capturas o informes.
El auto-registro de dispositivos se mantuvo desactivado:

\begin{lstlisting}
NGINX_TEMPLATE=https.conf.template
VEHICLESENSE_AUTO_REGISTER_DEVICES=false
\end{lstlisting}

\section{Primer arranque y migraciones}

El script \code{deploy.sh} valida los placeholders, comprueba el formato de la
contraseña de PostgreSQL, fuerza MQTT/TLS en el puerto 8883, valida Compose,
construye las imágenes y converge los servicios con \code{--wait}. En el primer
arranque se utilizó temporalmente \code{http.conf.template}; esto permite que
Nginx sirva el desafío ACME antes de disponer del certificado.

El entrypoint de FastAPI ejecutó:

\begin{lstlisting}
alembic upgrade head
\end{lstlisting}

La revisión aplicada fue \code{cd5ddc949e1d}, que quedó en \code{head}. El
administrador bootstrap se creó correctamente y los cuatro servicios
\code{postgres}, \code{backend}, \code{frontend} y \code{nginx} quedaron
\status{healthy}.

\section{HTTPS con Let's Encrypt}

Con DNS propagado y HTTP disponible, Certbot se ejecutó como servicio temporal de
Compose en modo webroot. Se emitió un certificado para el dominio principal y
\code{www}. Luego se cambió \code{NGINX\_TEMPLATE} a
\code{https.conf.template}, se recreó Nginx y se comprobó la redirección de
\code{www}, el proxy de la API y el canal WebSocket.

Durante el procedimiento se detectó un incidente de configuración: una sesión SSH
con variables exportadas mediante \code{source .env} conservaba el valor anterior
de \code{NGINX\_TEMPLATE} (\code{http.conf.template}), aunque el archivo ya había
sido editado a HTTPS. Compose prioriza las variables exportadas sobre el archivo
indicado con \code{--env-file}. La solución fue cerrar la sesión, abrir una sesión
limpia y recrear Nginx usando explícitamente \code{--env-file .env}. Desde entonces, las
operaciones normales utilizan los scripts versionados y no dejan secretos
exportados permanentemente en la shell.

\section{Renovación y operación}

La renovación se programó con una entrada cron diaria a las 03:17. El script de
renovación ejecuta \code{certbot renew} y recarga Nginx. Las
comprobaciones periódicas incluyen:

\begin{lstlisting}
docker compose --env-file deploy/.env -f deploy/compose.production.yml ps
docker compose --env-file deploy/.env -f deploy/compose.production.yml logs --tail=100
curl --fail https://vehiclemonitorsense.me/health/live
curl --fail https://vehiclemonitorsense.me/health/ready
df -h
free -h
swapon --show
\end{lstlisting}

Las actualizaciones se realizan con \code{update.sh}. Este exige un árbol Git
limpio, acepta únicamente fast-forward, genera un backup y guarda la revisión
anterior antes de construir la nueva versión.

\chapter{Herramientas y decisiones tecnológicas}

\section{Por qué AWS EC2}

AWS EC2 fue elegido por cuatro razones principales:

\begin{enumerate}
  \item \textbf{Control operativo.} Permite administrar Ubuntu, Docker, Nginx,
        cron, almacenamiento y red para demostrar un despliegue reproducible.
  \item \textbf{Servicios integrados.} Elastic IP, EBS, Security Groups y
        VPC resuelven direccionamiento, almacenamiento y control de acceso
        sin introducir un proveedor adicional para la máquina.
  \item \textbf{Compatibilidad y flexibilidad.} La familia Graviton ARM64 ofrece
        una instancia pequeña para comenzar, mientras que el mismo Compose puede
        ejecutarse en AMD64 o escalar verticalmente si aumentan los recursos.
  \item \textbf{Continuidad del proyecto.} La cuenta AWS, el dominio y la
        experiencia previa con Linux/SSH permitieron administrar el sistema con
        herramientas conocidas y conservar la posibilidad de ampliar la
        infraestructura más adelante.
\end{enumerate}

La decisión implica trabajo adicional: el sistema operativo, los parches, los
backups, los certificados, la monitorización y la seguridad no quedan delegados a
una plataforma PaaS. Ese costo operativo se aceptó porque el objetivo del proyecto
incluye comprender el ciclo completo de despliegue.

\section{Comparación con alternativas}

\begin{table}[ht]
\centering
\small
\begin{tabularx}{\textwidth}{>{\bfseries}l X X}
\toprule
Alternativa & Ventaja principal & Motivo para no usarla como plataforma final \\
\midrule
AWS EC2 & Control completo de VM, red, almacenamiento y seguridad; arquitectura
  portable con Docker & Requiere administrar el servidor y sus costos \\
Render/Railway & Despliegue sencillo y menor administración del sistema & Menor
  control de red, almacenamiento y proceso MQTT persistente; costos y límites
  dependen del plan \\
Vercel/Netlify & Excelente entrega de frontend estático & No resuelven por sí solos
  FastAPI, PostgreSQL, MQTT y WebSocket persistente \\
VPS administrada & Puede ser económica y simple para una VM & Menor integración
  nativa con la cuenta AWS y menor continuidad con la infraestructura elegida \\
OCI Compute & Alternativa portable para Docker y máquinas ARM & Se decidió
  mantener AWS como proveedor de referencia del proyecto \\
\bottomrule
\end{tabularx}
\caption{Comparación cualitativa de plataformas para este proyecto.}
\end{table}

La comparación no afirma que AWS sea universalmente superior. Para una aplicación
pequeña cuyo único objetivo fuera servir contenido estático, una plataforma PaaS
podría reducir el trabajo. En VehicleSense, la combinación de MQTT, persistencia,
WebSocket y control de red justificó la elección de una VM administrable.

\section{Herramientas utilizadas}

\begin{table}[ht]
\centering
\small
\begin{tabularx}{\textwidth}{>{\bfseries}l X}
\toprule
Herramienta & Uso realizado \\
\midrule
EC2 Console & Lanzamiento de la instancia, selección ARM64, llave y Security Group \\
Elastic IP & Dirección pública estable para DNS \\
Namecheap & Registros A para el dominio raíz y \code{www} \\
SSH/OpenSSH & Administración remota segura de Ubuntu \\
Docker Engine/Compose & Construcción y ejecución reproducible de los servicios \\
Nginx & Proxy inverso, HTTPS, WebSocket, headers y redirecciones \\
Certbot/Let's Encrypt & Emisión y renovación automática de certificados TLS \\
HiveMQ Cloud & Broker MQTT externo con TLS y ACL \\
FastAPI/PostgreSQL/Alembic & API, ingesta, persistencia y migraciones \\
React/Vite & Compilación y entrega del dashboard \\
Git y scripts Bash & Versionado, despliegue, actualización y backup \\
curl/dig/jq & Diagnóstico de DNS, HTTP, JSON y health checks \\
\bottomrule
\end{tabularx}
\caption{Herramientas principales del despliegue.}
\end{table}

\chapter{Seguridad, mantenimiento y costos}

\section{Controles aplicados}

La seguridad se planteó por capas:

\begin{itemize}
  \item el navegador nunca recibe credenciales MQTT ni PostgreSQL;
  \item el broker utiliza una credencial separada para el backend y TLS en 8883;
  \item el Security Group no expone PostgreSQL, FastAPI ni el frontend interno;
  \item Nginx aplica HTTPS, HSTS, CSP, \code{nosniff}, bloqueo de frames y límite
        de solicitudes para el login;
  \item FastAPI valida roles, payloads, tópicos, identidades y duplicados;
  \item los secretos permanecen en \code{deploy/.env}, con permisos restrictivos,
        y no se incluyen en la imagen del frontend;
  \item el acceso SSH se restringe a la IP del administrador y la llave privada
        no se almacena en el repositorio.
\end{itemize}

Estos controles reducen la superficie de ataque, pero no convierten al prototipo en
un sistema de seguridad vehicular certificado. Persisten riesgos como JWT sin
revocación persistente, ausencia de MFA, falta de ACL por vehículo en la API, OTA
local sin firma y una sola instancia.

\section{Backups y recuperación}

El backup lógico se crea antes de una actualización y también puede ejecutarse de
forma manual. El archivo dump se acompaña de un checksum y debe copiarse fuera de
EC2 a almacenamiento cifrado. La restauración es una operación explícita y
destructiva que requiere \code{CONFIRM\_RESTORE=vehiclesense}; primero se detiene
FastAPI y luego se reinicia después de \code{pg\_restore}.

No se borra el volumen \code{postgres\_data} para resolver fallos de migración.
Antes de cambiar una revisión se revisa la compatibilidad entre el código y las
migraciones.

\section{Memoria, costos y continuidad}

La instancia \code{t4g.small} tiene 2 GiB de RAM. Se recomienda comprobar swap,
memoria y disco antes de builds; si se presenta OOM se puede aumentar verticalmente
la instancia sin modificar Compose. La Elastic IP y el volumen EBS generan costos
según la cuenta y la región, por lo que deben revisarse las alertas de presupuesto
y evitarse recursos no utilizados.

La continuidad 24/7 depende de que EC2 permanezca ejecutándose, el dominio conserve
los registros DNS, el certificado se renueve y HiveMQ esté disponible. Apagar la
Mac, cerrar SSH o cambiar de red doméstica no detiene la plataforma web.

\chapter{Validación y resultados}

\section{Validaciones ejecutadas}

La validación se realizó en dos niveles. Primero se comprobó localmente que las
imágenes y la configuración fueran reproducibles. Después se verificaron los
endpoints públicos y el estado reportado de EC2.

\begin{table}[ht]
\centering
\small
\begin{tabularx}{\textwidth}{>{\bfseries}l X l}
\toprule
Prueba & Resultado & Alcance \\
\midrule
Build Docker & Backend y frontend construidos para ARM64 & Automatizado/local \\
Compose & Configuración válida; solo Nginx publica 80/443 & Automatizado/local \\
Migraciones & \code{cd5ddc949e1d (head)} & PostgreSQL de producción \\
Servicios & PostgreSQL, backend, frontend y Nginx healthy & Producción reportada \\
HTTPS/DNS & Dominio principal 200; \code{www} redirige; HTTP fuerza HTTPS & Público, 31/08/2026 \\
Health & \code{/health/live}=200 y \code{/health/ready}=200 & Público, 31/08/2026 \\
MQTT & Backend conectado por TLS/8883 & Producción reportada y health ready \\
Frontend/API & Página y login disponibles detrás de Nginx & Local y público \\
Backup/restore & Dump, checksum y restauración ejecutados localmente & Automatizado/local \\
\bottomrule
\end{tabularx}
\caption{Matriz resumida de validación del despliegue.}
\end{table}

La respuesta final de \code{/health/ready} indicó que la base de datos estaba
disponible y que el cliente MQTT estaba conectado. Esto valida la cadena cloud
hasta sus dependencias inmediatas, pero no reemplaza una sesión prolongada con el
ESP32 ni una medición de precisión de sensores.

\section{Incidentes y lecciones aprendidas}

El principal incidente fue el uso de variables exportadas de una sesión SSH
antigua. La shell mantuvo la plantilla HTTP y provocó que Nginx no escuchara en
443 aunque el archivo \code{.env} ya indicara HTTPS. La solución fue reiniciar la
sesión y usar \code{docker compose --env-file}. La lección es no mezclar variables
persistentes de shell con configuración versionada y validar siempre la
configuración efectiva de Nginx con \code{nginx -T}.

También se comprobó que el problema inicial de SSH provenía de la red local que
bloqueaba o interfería con TCP/22. El cambio temporal al hotspot permitió acceder
a EC2 y ajustar la regla del Security Group a la nueva IP. La web no depende del
hotspot; solo la administración SSH puede depender de esa red.

\section{Resultado funcional}

El resultado final es una página web accesible desde Internet, independiente de la
Mac, con autenticación, frontend React, API FastAPI, WebSocket, PostgreSQL y
conexión MQTT/TLS. Nginx termina HTTPS y mantiene los servicios internos aislados.
El despliegue es reproducible desde el repositorio mediante \code{deploy.sh} y las
actualizaciones se realizan con backup previo mediante \code{update.sh}.

\chapter{Procedimiento operativo}

\section{Comprobación diaria}

La secuencia mínima para revisar la plataforma es:

\begin{enumerate}
  \item entrar por SSH con la llave y la IP permitida;
  \item verificar \code{docker compose ps};
  \item revisar logs recientes de backend y Nginx;
  \item consultar \code{/health/live} y \code{/health/ready};
  \item revisar disco, memoria, swap y el cron de renovación;
  \item confirmar que no existan puertos internos publicados.
\end{enumerate}

\section{Actualización controlada}

Las actualizaciones no se hacen con un \code{git pull} aislado. Se ejecuta:

\begin{lstlisting}
cd /opt/vehiclesense
./deploy/scripts/update.sh main
\end{lstlisting}

El script exige un árbol limpio, toma un backup, registra la revisión previa,
acepta fast-forward y vuelve a construir el stack. Si el build falla, no ejecuta
rollback de la base de datos automáticamente: la decisión se toma revisando
migraciones y el backup.

\section{Registro de dispositivos}

Antes de aceptar telemetría real se registran en el backend los identificadores
exactos de vehículo y dispositivo. Deben coincidir con el firmware, el tópico MQTT
y las ACL de HiveMQ. Mantener \code{VEHICLESENSE\_AUTO\_REGISTER\_DEVICES=false}
evita que un dispositivo desconocido sea aceptado automáticamente.

\chapter{Conclusiones}

El despliegue de VehicleSense en AWS EC2 permitió convertir una aplicación local
en una plataforma web accesible permanentemente. La solución integra una VM ARM64,
Docker Compose, Nginx, FastAPI, PostgreSQL, HiveMQ Cloud y React sin exponer los
servicios internos. La configuración de DNS, Elastic IP, Security Group y
Let's Encrypt completó la publicación del dominio con HTTPS.

AWS EC2 fue una decisión coherente con el objetivo técnico y académico: ofreció
control directo de Linux y de la red, compatibilidad con Docker, herramientas de
operación conocidas y una ruta de crecimiento vertical. La contrapartida es que
el mantenimiento queda bajo responsabilidad del proyecto, especialmente backups,
parches, costos, monitorización y recuperación.

El sistema queda operativo para la siguiente fase: registrar los IDs reales,
conectar el ESP32, comprobar ingestión y WebSocket con telemetría física, y
documentar una sesión prolongada. Estas tareas amplían la validación del sistema,
pero no requieren volver a aprovisionar EC2, DNS, TLS ni HiveMQ.

\appendix
\chapter{Checklist de primer despliegue}

\begin{longtable}{p{0.06\textwidth} p{0.86\textwidth}}
\toprule
\# & Verificación \\
\midrule
1 & Instancia Ubuntu 24.04 ARM64/AMD64 creada y Elastic IP asociada. \\
2 & Security Group con SSH restringido y HTTP/HTTPS públicos; sin 5432, 8000, 8080 ni 8883 de entrada. \\
3 & DNS A de \code{@} y \code{www} resuelve a la Elastic IP. \\
4 & Docker Engine y Compose instalados; usuario reconectado al grupo Docker. \\
5 & \code{deploy/.env} creado con permisos 600 y sin placeholders. \\
6 & Credenciales de PostgreSQL, JWT, administrador y HiveMQ separadas. \\
7 & \code{NGINX\_TEMPLATE=http.conf.template} usado solo para el primer desafío ACME. \\
8 & \code{deploy.sh} ejecutado; cuatro servicios healthy y Alembic en head. \\
9 & Certbot emitió certificado para dominio raíz y \code{www}. \\
10 & \code{NGINX\_TEMPLATE=https.conf.template}; HTTP redirige a HTTPS. \\
11 & Renovación cron configurada una sola vez y verificada con \code{crontab -l}. \\
12 & Bootstrap admin vaciado después de confirmar el acceso. \\
13 & Backup con checksum copiado a almacenamiento externo cifrado. \\
14 & IDs de vehículo/dispositivo registrados y telemetría física validada. \\
\bottomrule
\end{longtable}

\renewcommand{\bibname}{Referencias}

\begin{thebibliography}{9}
\bibitem{aws-sg}
Amazon Web Services, ``Security groups for your EC2 instances'',
\url{https://docs.aws.amazon.com/AWSEC2/latest/UserGuide/ec2-security-groups.html}.

\bibitem{aws-eip}
Amazon Web Services, ``Elastic IP addresses'',
\url{https://docs.aws.amazon.com/AWSEC2/latest/UserGuide/elastic-ip-addresses-eip.html}.

\bibitem{docker-compose}
Docker, ``Docker Compose documentation'',
\url{https://docs.docker.com/compose/}.

\bibitem{certbot}
Certbot, ``Using Certbot with webroot'',
\url{https://eff-certbot.readthedocs.io/en/stable/using.html}.

\bibitem{hivemq}
HiveMQ, ``HiveMQ Cloud documentation'',
\url{https://www.hivemq.com/docs/hivemq-cloud/}.

\bibitem{proyecto}
VehicleSense, ``Documentación interna del proyecto: arquitectura, seguridad,
despliegue y pruebas'', repositorio del proyecto, 2026.
\end{thebibliography}

\end{document}
